% Irodov Part III §3.6 original problems 305--322, source printed pp.139--141,
% PDF 148--150; corresponding source answers printed pp.308--309.
\electroinductionproblem{305}{139}
图3.91平面回路由边长 $a=20\,\mathrm{cm}$ 和 $b=10\,\mathrm{cm}$ 的两个方形回路组成，处于垂直纸面的变化磁场 $B=B_0\sin\omega t$，$B_0=10\,\mathrm{mT}$、$\omega=100\,\mathrm{s^{-1}}$。导线单位长度电阻为 $\rho=50\,\mathrm{m\Omega/m}$。忽略自感，求感生电流振幅。
\exsolution
按图3.91的绕线方向，两个方环的有向面积相减，有效磁通面积 $A_{\rm eff}=a^2-b^2$。两环总导线长为 $4(a+b)$，电阻 $R_{\rm loop}=4\rho(a+b)$。磁通量为 $\Phi=(a^2-b^2)B_0\sin\omega t$，故电动势振幅
$\mathcal E_m=\omega B_0(a^2-b^2)$。电流振幅
\[
\boxed{I_m=\frac{\omega B_0(a^2-b^2)}{4\rho(a+b)}
=\frac{\omega B_0(a-b)}{4\rho}
=0.50\,\mathrm A}.
\]
\excheck 要按图判断两方环有向面积相减，否则错用 $a^2+b^2$；分母的导线长度也是两个方环周长之和。

\electroinductionproblem{306}{139}
紧密绕制、匝数很多的单层平面螺旋，外圈半径为 $a$，总匝数为 $N$。均匀磁场垂直螺旋平面且 $B(t)=B_0\sin\omega t$，求感应电动势振幅。
\exsolution
密绕螺旋的每匝半径近似从零至 $a$ 均匀分布，每匝所围面积为 $\pi r^2$。匝数密度为 $\mathrm dN=N\,\mathrm dr/a$，总磁链
\[
\Psi=B(t)\int_0^a\pi r^2\frac Na\,\mathrm dr
=B(t)\frac{\pi Na^2}{3}.
\]
法拉第定律 $\mathcal E=-\dot\Psi$ 给
\[
\boxed{\mathcal E_m=\frac{\pi Na^2\omega B_0}{3}}.
\]
\excheck 不能直接用 $N\pi a^2$，因为内圈面积比外圈小；均方半径为 $a^2/3$。

\electroinductionproblem{307}{139}
U 形导轨位于磁场垂直纸面的区域，磁感应强度随时间以 $\dot B=0.10\,\mathrm{T/s}$ 均匀增加。长为 $\ell=20\,\mathrm{cm}$ 的活动杆从静止开始以 $w=10\,\mathrm{cm/s^2}$ 加速滑动，且 $t=0$ 时回路面积与磁场均为零。求 $t=2.0\,\mathrm s$ 时感应电动势。
\exsolution
活动杆位移 $x=wt^2/2$，回路面积 $S=\ell wt^2/2$。磁场 $B=\dot B\,t$，所以
$\Phi=BS=\dot B\,\ell w t^3/2$。由法拉第定律
\[
\boxed{|\mathcal E(t)|=\left|\frac{\mathrm d\Phi}{\mathrm dt}\right|
=\frac32\dot B\,\ell w t^2}.
\]
代入 $t=2$ s：
$\boxed{|\mathcal E|=0.012\,\mathrm V=12\,\mathrm{mV}}$。
\excheck 磁通量对时间的导数既包含 $\dot B S$，也包含 $B\dot S$；只计算其中一项会错系数。

\electroinductionproblem{308}{139}
无限长密绕螺线管半径为 $a$，单位长度匝数 $n$，其电流按恒定速率 $\dot I$ 变化。求涡旋感生电场大小随离轴距离 $r$ 的关系并画出定性变化。
\exsolution
理想螺线管内部 $B=\mu_0nI$，外部磁场近似为零。取半径 $r$ 的同轴圆周为法拉第环路。对 $r<a$：
\[
2\pi rE_\varphi=\left|\frac{\mathrm d}{\mathrm dt}(B\pi r^2)\right|
=\pi r^2\mu_0n|\dot I|,
\quad \boxed{E_\varphi=\frac12\mu_0n|\dot I|r}.
\]
对 $r>a$，磁通只通过管截面：
\[
2\pi rE_\varphi=\pi a^2\mu_0n|\dot I|,
\quad\boxed{E_\varphi=\frac{\mu_0n|\dot I|a^2}{2r}}.
\]
\excheck $E_\varphi$ 在 $r=a$ 连续；内部按 $r$ 上升，外部按 $1/r$ 下降。此处电场非保守，环流非零。

\electroinductionproblem{309}{139}
很长螺线管直径 $d=5.0\,\mathrm{cm}$，单位长度匝数 $n=20/\mathrm{cm}$，线圈电流以 $\dot I=100\,\mathrm{A/s}$ 增加。在螺线管上紧套一圈铜线，其截面积 $S=1.0\,\mathrm{mm^2}$，忽略该圈自感，求感生电流。
\exsolution
线圈内部磁场变化率
$\dot B=\mu_0n\dot I$。单匝围面积 $\pi d^2/4$，
其感生电动势
$\mathcal E=(\pi d^2/4)\mu_0n|\dot I|$。
铜圈长度近似 $\pi d$，电阻
$R=\rho_{\rm Cu}\pi d/S$，其中 $\rho_{\rm Cu}$ 为电阻率。于是
\[
\boxed{I_{\rm ind}=\frac{\mu_0ndS|\dot I|}{4\rho_{\rm Cu}}}.
\]
取铜电阻率 $1.7\times10^{-8}\,\Omega\cdot\mathrm m$，
数值 $\boxed{I_{\rm ind}\simeq2.0\,\mathrm{mA}}$。
\excheck 题目并未明给铜电阻率，数值需声明采用常温近似；导线截面积 $S$ 不等于螺线管横截面积。

\electroinductionproblem{310}{139}
长螺线管半径为 $a$，其外紧套一圈细绝缘导线闭环，两段半圆形导线中，一段电阻是另一段的 $\eta$ 倍。螺线管磁场按 $B(t)=bt$ 变化，求环中建立的静电场大小。
\exsolution
电磁感应产生的环向非保守电场
$E_{\rm ind}=-ab/2$（选择有向环路）。
一圈总电动势为 $|\mathcal E|=\pi a^2|b|$。
两半环长度相同但电阻分别为 $R_0,\eta R_0$，
因此各段\emph{总驱动力场}大小按其分担的电压成比例，
分别为
$E_{{\rm tot},1}=ab/(1+\eta)$、
$E_{{\rm tot},2}=ab\eta/(1+\eta)$。
总驱动力场是非保守感生电场与由面电荷建立的保守静电场之和；故静电场在两半环的大小相同、方向相反，
\[
\boxed{E_{\rm es}=\frac{|ab|}{2}
\frac{|\eta-1|}{\eta+1}}.
\]
\excheck $\eta=1$ 时两段相同，无需额外静电场重新分配感生驱动力，故 $E_{\rm es}=0$；必须区分静电场与总驱动力场。

\electroinductionproblem{311}{139}
质量为 $m$、总电荷 $q$ 的薄绝缘圆环可绕自身轴自由转动。初始静止且无磁场，此后垂直环面的磁场从零按任意 $B(t)$ 增强。求角速度 $\omega(t)$。
\exsolution
设圆环半径为 $a$。法拉第定律给环周涡旋场
$E_\varphi(a)=-a\dot B/2$。圆环受电力矩
$\tau=qE_\varphi(a)a=-qa^2\dot B/2$。
薄圆环转动惯量 $J=ma^2$，由
$J\dot\omega=\tau$，
\[
\dot\omega=-\frac q{2m}\dot B.
\]
初始 $\omega(0)=B(0)=0$，积分得到
\[
\boxed{\omega(t)=-\frac q{2m}B(t)}.
\]
\excheck 半径消去；角速度符号依圆环电荷与磁场正方向而定。

\electroinductionproblem{312}{140}
长为 $\ell$、截面半径为 $b$ 的螺线管内部，同轴放一电阻为 $r$、半径为 $a$ 的细金属环。螺线管突然接电压 $V$ 的直流电源，电源及螺线管总电阻为 $R$。忽略小环自感，求小环每单位长度所受径向磁力的最大值。
\exsolution
设螺线管共有 $N$ 匝，自感
$L_s=\mu_0N^2\pi b^2/\ell$。主电流
$I_s(t)=(V/R)(1-e^{-tR/L_s})$，
主磁场
$B(t)=\mu_0NI_s(t)/\ell$。
环内磁通 $\Phi=\pi a^2B$，
感应电流
$|I_r(t)|=(\pi a^2/r)\dot B
=(a^2V/(rNb^2))e^{-tR/L_s}$。
环的径向安培力单位长度模
$f(t)=B(t)|I_r(t)|$，故
\[
f(t)=\frac{\mu_0a^2V^2}{rR\ell b^2}
(1-e^{-u})e^{-u},\quad u=tR/L_s.
\]
$x(1-x)$ 于 $x=e^{-u}=1/2$ 最大值为 $1/4$，所以
\[
\boxed{f_{\max}=\frac{\mu_0a^2V^2}{4rR\ell b^2}}.
\]
\excheck 螺线管匝数 $N$ 最终消去；最大值在 $t=(L_s/R)\ln2$，而不是电流刚建立或最终稳定时。

\electroinductionproblem{313}{140}
固定闭合导线回路的电阻为 $R$。在 $0\le t\le\tau$ 内穿过回路的磁通量
$\Phi(t)=at(\tau-t)$，忽略自感，求这段时间内的焦耳热。
\exsolution
由法拉第定律
$\mathcal E=-\dot\Phi=-a(\tau-2t)$，
$I=\mathcal E/R$。
故
\[
Q=\int_0^\tau I^2R\,\mathrm dt
=\frac{a^2}R\int_0^\tau(\tau-2t)^2\,\mathrm dt.
\]
积分
$\int_0^\tau(\tau-2t)^2\,\mathrm dt=\tau^3/3$，因此
\[
\boxed{Q=\frac{a^2\tau^3}{3R}}.
\]
\excheck 虽然首末磁通均为零，过程中的感应电流正负变向，但焦耳热是正值。

\electroinductionproblem{314}{140}
长螺线管中央同轴放置一个矩形截面导体圆环，内外半径为 $a,b$，轴向厚度为 $h$，材料电阻率为 $\rho$。螺线管磁场按 $B(t)=\beta t$ 变化。忽略环自身电感，求总环向感应电流。
\exsolution
对半径为 $r$ 的圆周，非保守感生场
$E_\varphi(r)=-r\beta/2$。
材料的环向电流密度模为
$j_\varphi(r)=|E_\varphi|/\rho=|\beta|r/(2\rho)$。
各半径薄层并联通过径向截面面积
$\mathrm dS=h\,\mathrm dr$，
总电流
\[
I=\int_a^bj_\varphi(r)h\,\mathrm dr
=\frac{|\beta|h}{2\rho}\frac{b^2-a^2}{2}.
\]
因此
\[
\boxed{|I|=\frac{|\beta|h(b^2-a^2)}{4\rho}}.
\]
\excheck 本题环是具有有限径向厚度的\emph{实体导电圆环}，不能当作单根半径平均值的细导线。

\electroinductionproblem{315}{140}
要制作长度为 $\ell_0=100\,\mathrm{cm}$、自感为 $L=1.0\,\mathrm{mH}$ 的长直细螺线管，直径远小于长度，问需用多长细导线？
\exsolution
令总匝数为 $N$、半径为 $R$，导线总长近似为 $s=2\pi RN$。长螺线管自感
\[
L=\frac{\mu_0N^2\pi R^2}{\ell_0}
=\frac{\mu_0s^2}{4\pi\ell_0}.
\]
解得
\[
\boxed{s=\sqrt{\frac{4\pi\ell_0L}{\mu_0}}
\simeq1.0\times10^2\,\mathrm m}.
\]
\excheck $4\pi/\mu_0=10^7\,\mathrm{H^{-1}m^{-1}}$，
$\ell_0L=10^{-3}\,\mathrm{H\,m}$，故根号下为 $10^4\,\mathrm{m^2}$，得 $100$m；原书答案若印 $1.10\,\mathrm{km}$，需重新核对公式的符号与数值，不能在译稿中掩盖这一矛盾。

\electroinductionproblem{316}{140}
一个长直螺线管长度为 $\ell$，绕线全部由质量 $m$ 的铜丝制成，铜丝总电阻为 $R$，且螺线管直径远小于长度。求自感 $L$；铜电阻率为 $\rho$，材料密度为 $\rho_m$。
\exsolution
设铜丝长度为 $s$、截面积为 $A$，
质量 $m=\rho_mAs$，
电阻 $R=\rho s/A$。相乘：
$mR=\rho_m\rho s^2$，所以
$s^2=mR/(\rho_m\rho)$。
长直单层密绕螺线管自感可写成
$L=\mu_0s^2/(4\pi\ell)$，于是
\[
\boxed{L=\frac{\mu_0mR}{4\pi\ell\rho\rho_m}}.
\]
\excheck 电阻率 $\rho$ 与材料\emph{质量密度} $\rho_m$ 量纲不同，不可混用。

\electroinductionproblem{317}{140}
自感为 $L=300\,\mathrm{mH}$、电阻为 $R=140\,\mathrm{m\Omega}$ 的线圈接直流恒压源，求电流达到稳态值的 $\eta=50\%$ 所需时间。
\exsolution
串联 RL 回路方程
$L\dot I+RI=\mathcal E$，
初始电流零时
$I(t)=I_\infty(1-e^{-Rt/L})$，
其中 $I_\infty=\mathcal E/R$。
令 $I/I_\infty=\eta$，
\[
\boxed{t=-\frac LR\ln(1-\eta)}.
\]
代入 $L=0.300$H、$R=0.140\,\Omega$、$\eta=0.50$，
$\boxed{t\simeq1.5\,\mathrm s}$。
\excheck 线圈电阻是 $140$mΩ 而非 $140$Ω，否则时间差千倍。

\electroinductionproblem{318}{140}
长 $\ell=1.0\,\mathrm m$ 的直螺线管由质量 $m=1.0\,\mathrm{kg}$ 的单层铜线绕成，直径远小于长度，求 RL 时间常数 $\tau=L/R$。
\exsolution
由第316题的结果
$L=\mu_0mR/(4\pi\ell\rho\rho_m)$，
故
\[
\boxed{\tau=\frac LR=\frac{\mu_0m}{4\pi\ell\rho\rho_m}}.
\]
取铜电阻率
$\rho\simeq1.7\times10^{-8}\,\Omega\cdot\mathrm m$、
密度 $\rho_m\simeq8.9\times10^3\,\mathrm{kg/m^3}$，得到
\[
\boxed{\tau\simeq6.6\times10^{-4}\,\mathrm s
\approx0.7\,\mathrm{ms}}.
\]
\excheck 与线圈的绕线半径和匝数无关是长直、紧密单层绕线近似的结果。

\electroinductionproblem{319}{140}
同轴双导体电缆由两个无限长薄壁圆柱形金属筒构成，外筒半径是内筒的 $\eta=3.6$ 倍，筒间磁导率为 $\mu_0$。求每单位长度的电感。
\exsolution
内筒与外筒通等量反向电流 $I$，仅筒间存在磁场
$B(r)=\mu_0I/(2\pi r)$。
单位长度储磁能为
\[
w'_m=\int_a^b\frac{B(r)^2}{2\mu_0}2\pi r\,\mathrm dr
=\frac{\mu_0I^2}{4\pi}\ln\frac ba.
\]
用 $w'_m=\tfrac12L'I^2$：
\[
\boxed{L'=\frac{\mu_0}{2\pi}\ln\eta}.
\]
$\eta=3.6$ 时
$\boxed{L'\simeq2.6\times10^{-7}\,\mathrm{H/m}
=0.26\,\mathrm{\mu H/m}}$。
\excheck 薄壁导体内部磁场忽略，故没有内电感修正。

\electroinductionproblem{320}{140}
截面为边长 $a$ 的正方形、内半径为 $b$ 的环形螺线管共有 $N$ 匝，介质相对磁导率为 $\mu$，求自感。
\exsolution
环向磁场由安培定律给
$B(r)=\mu\mu_0NI/(2\pi r)$，
$r\in[b,b+a]$。
每匝截面磁通
\[
\Phi=\int_b^{b+a}B(r)\,a\,\mathrm dr
=\frac{\mu\mu_0NIa}{2\pi}\ln\frac{b+a}{b}.
\]
总磁链 $N\Phi=LI$，所以
\[
\boxed{L=\frac{\mu\mu_0N^2a}{2\pi}
\ln\left(1+\frac ab\right)}.
\]
\excheck $a\ll b$ 时近似
$L\simeq\mu\mu_0N^2a^2/(2\pi b)$，与均匀环形磁路近似一致。

\electroinductionproblem{321}{141}
图3.92两条宽 $b$、相距 $h$ 的平行宽薄金属带，$b/h=50$，两条导体中电流等量反向。求每单位长度电感。
\exsolution
忽略边缘效应，单位宽度表面电流为 $K=I/b$，
两带之间磁场近似均匀
$B=\mu_0K=\mu_0I/b$，外部磁场近似零。
单位长度磁场储能
$w'_m=(B^2/2\mu_0)(bh)
=\mu_0I^2h/(2b)$。
比较 $w'_m=L'I^2/2$：
\[
\boxed{L'=\frac{\mu_0h}{b}}.
\]
代 $b/h=50$：
$\boxed{L'\simeq25\,\mathrm{nH/m}}$。
\excheck 宽带近似要求 $h\ll b$；边缘漏磁只贡献较小修正。

\electroinductionproblem{322}{141}
两根无限长平行圆导线，轴间距离为每根线半径的 $\eta\gg1$ 倍，两线通等量反向电流。忽略导线内部磁场，求双线系统的单位长度电感。
\exsolution
对轴间距 $d$、半径 $a$ 的远离双线，
外部磁通/能量计算可用两条无限线电流的叠加；每单位长度电感准确形式为
$L'=(\mu_0/\pi)\operatorname{arcosh}(d/2a)$。
在 $d/a=\eta\gg1$ 时
$\operatorname{arcosh}(\eta/2)\simeq\ln\eta$，
因而
\[
\boxed{L'\simeq\frac{\mu_0}{\pi}\ln\eta}.
\]
\excheck 本题排除导线内部场能，因此不加圆导线的内电感 $\mu_0/(8\pi)$ 一类修正。
